\begin{pastexam}[label={px:ntust-112-ice-04}]{臺科大資工所 112 學年度 第 4 題}
（問答）

(5\%) Alice tries to insert a set of keys $[17,88,62,30,28,77,90,45,51]$ sequentially into a fresh
hash table with size 3 based on the linear probing strategy $g(x,i)=(h(x)+i)\bmod M$, where $M$ is
the size of the hash table and $h(x)=x\bmod M$. If the overflow condition happens, rehash will be
performed by doubling the table size and computing the new hash value for the elements from the
beginning in the old table. Following that, the new table continues to be filled with uninserted
elements. Please write down the final result of the hash table.
\end{pastexam}

\begin{pastsoln}
答案：最後的表大小 $M = 12$：
\begin{center}
\small
\begin{tabular}{l*{12}{c}}
\toprule
位置 & 0 & 1 & 2 & 3 & 4 & 5 & 6 & 7 & 8 & 9 & 10 & 11 \\
\midrule
鍵 & -- & -- & 62 & 51 & 28 & 77 & 30 & 88 & 17 & 90 & 45 & -- \\
\bottomrule
\end{tabular}
\end{center}

「overflow」讀成「要插入時表已經滿了，線性探測找不到空位」。 這時表的大小加倍，依舊表的位置
$0, 1, 2, \ldots$ 的順序把舊的鍵用新的 $M$ 重新插入，再繼續插入還沒放的鍵（資結書第~12.3~節的線性
探測、第~12.5~節的加倍再雜湊）。

\textbf{$M = 3$。}
\begin{itemize}
\item $17 \bmod 3 = 2$，放在 $2$。
\item $88 \bmod 3 = 1$，放在 $1$。
\item $62 \bmod 3 = 2$，$2$ 被占，試 $(2 + 1) \bmod 3 = 0$，放在 $0$。 表：$[62, 88, 17]$，滿了。
\item $30$：表滿，overflow，加倍成 $M = 6$。
\end{itemize}

\textbf{$M = 6$。} 先依舊表的順序重新插入 $62, 88, 17$：$62 \bmod 6 = 2$；$88 \bmod 6 = 4$；
$17 \bmod 6 = 5$。 表：$[-, -, 62, -, 88, 17]$。 接著繼續：
\begin{itemize}
\item $30 \bmod 6 = 0$，放在 $0$。
\item $28 \bmod 6 = 4$，$4$、$5$、$0$ 被占，放在 $1$。
\item $77 \bmod 6 = 5$，$5$、$0$、$1$、$2$ 被占，放在 $3$。 表：$[30, 28, 62, 77, 88, 17]$，滿了。
\item $90$：表滿，overflow，加倍成 $M = 12$。
\end{itemize}

\textbf{$M = 12$。} 依舊表的順序重新插入 $30, 28, 62, 77, 88, 17$：
\begin{itemize}
\item $30 \bmod 12 = 6$，放在 $6$；$28 \bmod 12 = 4$，放在 $4$；$62 \bmod 12 = 2$，放在 $2$；
$77 \bmod 12 = 5$，放在 $5$。
\item $88 \bmod 12 = 4$，$4$、$5$、$6$ 被占，放在 $7$。
\item $17 \bmod 12 = 5$，$5$、$6$、$7$ 被占，放在 $8$。
\end{itemize}
再繼續插入剩下的鍵：
\begin{itemize}
\item $90 \bmod 12 = 6$，$6$、$7$、$8$ 被占，放在 $9$。
\item $45 \bmod 12 = 9$，$9$ 被占，放在 $10$。
\item $51 \bmod 12 = 3$，放在 $3$。
\end{itemize}
九個鍵都放完，得到上面的表。 注意重新插入的順序會影響結果：例如 $88$ 和 $28$ 在 $M = 12$ 時都雜湊到
$4$，照舊表的順序是 $28$（舊位置 $1$）先放、$88$（舊位置 $4$）後放，所以 $28$ 在 $4$、$88$ 被擠到 $7$。
\end{pastsoln}
